% ============================================================
\chapter{Fonctions de Variables Al\'eatoires}
\label{ch:fonctions_va}
% ============================================================

En sciences et en ing\'enierie, on mesure rarement la quantit\'e qui nous int\'eresse directement. Un physicien mesure un angle et en d\'eduit une vitesse~; un financier observe un rendement logarithmique et en d\'eduit un prix~; un ing\'enieur mesure une tension et en d\'eduit une puissance. \`A chaque fois, on applique une transformation $g$ \`a une variable al\'eatoire $X$ pour obtenir une nouvelle variable al\'eatoire $Y = g(X)$. Mais quelle est la loi de $Y$~? La question, d'apparence innocente, conduit \`a des techniques riches et vari\'ees~: formule du changement de variable, m\'ethode de la fonction de r\'epartition, jacobiens pour les transformations multidimensionnelles. Ma\^itriser ces outils, c'est passer du statut d'observateur passif de l'al\'ea \`a celui d'architecte capable de manipuler les distributions \`a volont\'e.

\begin{intuition}
Lorsqu'on applique une transformation \`a une variable al\'eatoire, on obtient une
nouvelle variable al\'eatoire dont la loi peut \^etre tr\`es diff\'erente. Ce chapitre
d\'eveloppe les m\'ethodes permettant de d\'eterminer la loi de $Y = g(X)$ \`a partir
de celle de $X$.
\end{intuition}

% ------------------------------------------------------------
\section{Cas discret}
% ------------------------------------------------------------

\begin{proposition}[Loi de $g(X)$ --- cas discret]
Si $X$ est une variable discr\`ete prenant les valeurs $x_1, x_2, \ldots$ et
$g : \R \to \R$, alors $Y = g(X)$ est discr\`ete et~:
\[
    \PP(Y = y) = \sum_{k : g(x_k) = y} \PP(X = x_k).
\]
\end{proposition}

\begin{exemple}
Si $X \sim \text{Bin}(n, 1/2)$ et $Y = 2X - n$, alors $Y$ prend les valeurs
$-n, -n+2, \ldots, n-2, n$ et~:
\[
    \PP(Y = 2k-n) = \binom{n}{k} \left(\frac{1}{2}\right)^n, \quad k = 0, 1, \ldots, n.
\]
\end{exemple}

% ------------------------------------------------------------
\section{M\'ethode de la fonction de r\'epartition}
% ------------------------------------------------------------

\begin{theoreme}[M\'ethode de la CDF]
\label{thm:methode_cdf}
Pour trouver la loi de $Y = g(X)$~:
\begin{enumerate}
    \item Calculer $F_Y(y) = \PP(Y \leq y) = \PP(g(X) \leq y)$.
    \item D\'eriver~: $f_Y(y) = F_Y'(y)$.
\end{enumerate}
\end{theoreme}

\begin{exemple}[$Y = X^2$ avec $X \sim \mathcal{N}(0,1)$]
Pour $y > 0$~:
\begin{align*}
    F_Y(y) &= \PP(X^2 \leq y) = \PP(-\sqrt{y} \leq X \leq \sqrt{y})
    = \Phi(\sqrt{y}) - \Phi(-\sqrt{y}) = 2\Phi(\sqrt{y}) - 1.
\end{align*}
En d\'erivant~:
\[
    f_Y(y) = 2 \cdot \frac{1}{\sqrt{2\pi}} e^{-y/2} \cdot \frac{1}{2\sqrt{y}}
    = \frac{1}{\sqrt{2\pi}}\, y^{-1/2}\, e^{-y/2}, \quad y > 0.
\]
C'est la densit\'e d'une loi $\Gamma(1/2, 1/2) = \chi^2(1)$.
\end{exemple}

\begin{exemple}[$Y = e^X$ avec $X \sim \mathcal{N}(\mu, \sigma^2)$]
Pour $y > 0$~:
\[
    F_Y(y) = \PP(e^X \leq y) = \PP(X \leq \ln y) = \Phi\!\left(\frac{\ln y - \mu}{\sigma}\right).
\]
En d\'erivant~:
\[
    f_Y(y) = \frac{1}{y\sigma\sqrt{2\pi}}
    \exp\!\left(-\frac{(\ln y - \mu)^2}{2\sigma^2}\right), \quad y > 0.
\]
On retrouve la densit\'e de la loi log-normale.
\end{exemple}

% ------------------------------------------------------------
\section{M\'ethode du changement de variable}
% ------------------------------------------------------------

\begin{theoreme}[Changement de variable --- cas monotone]
\label{thm:changement_var}
Soit $X$ une variable continue de densit\'e $f_X$ et $g$ une fonction
\textbf{strictement monotone} et de classe $C^1$ sur le support de $X$, avec
$g^{-1}$ de classe $C^1$. Alors $Y = g(X)$ a pour densit\'e~:
\[
    f_Y(y) = f_X\bigl(g^{-1}(y)\bigr) \cdot \abs{(g^{-1})'(y)}.
\]
\end{theoreme}

\begin{proof}
Supposons $g$ strictement croissante (le cas d\'ecroissant est analogue).
\begin{align*}
    F_Y(y) = \PP(g(X) \leq y) = \PP(X \leq g^{-1}(y)) = F_X(g^{-1}(y)).
\end{align*}
En d\'erivant par la r\`egle de la cha\^ine~:
$f_Y(y) = f_X(g^{-1}(y)) \cdot (g^{-1})'(y)$.
Comme $g$ est croissante, $(g^{-1})' > 0$, d'o\`u la valeur absolue est superflue.
Dans le cas d\'ecroissant, on obtient un signe moins qui est absorb\'e par $\abs{\cdot}$.
\end{proof}

\begin{exemple}[Transformation affine]
Si $Y = aX + b$ avec $a \neq 0$~:
$g^{-1}(y) = (y-b)/a$, $(g^{-1})'(y) = 1/a$. Donc~:
\[
    f_Y(y) = \frac{1}{\abs{a}}\, f_X\!\left(\frac{y-b}{a}\right).
\]
En particulier, si $X \sim \mathcal{N}(\mu, \sigma^2)$, alors
$Z = (X-\mu)/\sigma \sim \mathcal{N}(0,1)$.
\end{exemple}

% ------------------------------------------------------------
\section{Cas non monotone}
% ------------------------------------------------------------

\begin{theoreme}[Changement de variable --- cas g\'en\'eral]
Si $g$ est de classe $C^1$ et le support de $X$ peut \^etre d\'ecompos\'e en un
nombre fini d'intervalles $I_1, \ldots, I_m$ sur chacun desquels $g$ est
strictement monotone, avec inverses $h_1, \ldots, h_m$, alors~:
\[
    f_Y(y) = \sum_{j=1}^{m} f_X(h_j(y)) \cdot \abs{h_j'(y)}\, \mathbf{1}_{g(I_j)}(y).
\]
\end{theoreme}

\begin{exemple}[$Y = X^2$ avec $X$ de densit\'e $f_X$ sym\'etrique]
$g(x) = x^2$ est d\'ecroissante sur $(-\infty, 0)$ et croissante sur $(0, +\infty)$.
Les deux inverses sont $h_1(y) = -\sqrt{y}$ et $h_2(y) = \sqrt{y}$, avec
$\abs{h_j'(y)} = 1/(2\sqrt{y})$. Donc~:
\[
    f_Y(y) = \frac{f_X(\sqrt{y}) + f_X(-\sqrt{y})}{2\sqrt{y}}, \quad y > 0.
\]
\end{exemple}

% ------------------------------------------------------------
\section{Fonction g\'en\'eratrice de moments}
% ------------------------------------------------------------

\begin{definition}[Fonction g\'en\'eratrice de moments (MGF)]
La \textbf{fonction g\'en\'eratrice de moments} de $X$ est~:
\[
    M_X(t) = \E[e^{tX}],
\]
d\'efinie pour les $t$ dans un voisinage de 0 o\`u cette esp\'erance existe.
\end{definition}

\begin{proposition}[Propri\'et\'es de la MGF]
\begin{enumerate}[label=(\roman*)]
    \item $M_X(0) = 1$.
    \item $M_X^{(n)}(0) = \E[X^n]$ (les moments s'obtiennent par d\'erivation).
    \item Si $M_X(t) = M_Y(t)$ pour tout $t$ dans un voisinage de 0, alors
          $X$ et $Y$ ont m\^eme loi.
    \item Si $X \indep Y$, alors $M_{X+Y}(t) = M_X(t) \cdot M_Y(t)$.
\end{enumerate}
\end{proposition}

\begin{exemple}[MGF de la loi normale]
Si $X \sim \mathcal{N}(\mu, \sigma^2)$~:
\[
    M_X(t) = \exp\!\left(\mu t + \frac{\sigma^2 t^2}{2}\right).
\]
On retrouve $M_X'(0) = \mu = \E[X]$ et
$M_X''(0) = \sigma^2 + \mu^2 = \E[X^2]$.
\end{exemple}

% ------------------------------------------------------------
\section{Transformation de probabilit\'e int\'egrale}
% ------------------------------------------------------------

\begin{theoreme}[Transformation de probabilit\'e int\'egrale]
Si $X$ est une variable continue de CDF $F_X$ strictement croissante, alors
$U = F_X(X) \sim \mathcal{U}([0,1])$.
\end{theoreme}

\begin{proof}
Pour $u \in (0,1)$~:
$\PP(U \leq u) = \PP(F_X(X) \leq u) = \PP(X \leq F_X^{-1}(u)) = F_X(F_X^{-1}(u)) = u$.
\end{proof}

\begin{remarque}
Ce r\'esultat est la base th\'eorique de la m\'ethode d'inversion pour la simulation
(chapitre~4) et des tests d'ad\'equation (p-values).
\end{remarque}

% ------------------------------------------------------------
\section{Exercices}
% ------------------------------------------------------------

\begin{exercice}
Soit $X \sim \mathcal{U}([0,1])$ et $Y = -\ln X$. D\'eterminer la loi de $Y$.
\end{exercice}

\begin{exercice}
Soit $X \sim \text{Exp}(1)$. Calculer la densit\'e de $Y = \sqrt{X}$.
\end{exercice}

\begin{exercice}
Soit $X \sim \mathcal{N}(0,1)$. Calculer la densit\'e de $Y = \abs{X}$
(loi normale repli\'ee).
\end{exercice}

\begin{exercice}
Soit $X$ de densit\'e $f_X(x) = 2x\, \mathbf{1}_{[0,1]}(x)$ et $Y = 1 - X^2$.
D\'eterminer la densit\'e de $Y$.
\end{exercice}

\begin{exercice}
Calculer la MGF de la loi $\text{Exp}(\lambda)$ et en d\'eduire $\E[X]$, $\E[X^2]$,
$\Var(X)$.
\end{exercice}

\begin{exercice}
Soit $X \sim \text{Cauchy}(0,1)$ de densit\'e $f(x) = 1/(\pi(1+x^2))$.
Montrer que la MGF de $X$ n'existe pas (i.e.\ $\E[e^{tX}] = \infty$ pour
tout $t \neq 0$).
\end{exercice}

\begin{exercice}
Soit $U_1, U_2$ i.i.d.\ $\mathcal{U}([0,1])$. Par la m\'ethode de Box--Muller,
on pose~:
$Z_1 = \sqrt{-2\ln U_1} \cos(2\pi U_2)$ et
$Z_2 = \sqrt{-2\ln U_1} \sin(2\pi U_2)$.
Montrer que $Z_1, Z_2$ sont i.i.d.\ $\mathcal{N}(0,1)$.
\end{exercice}

\begin{exercice}
En utilisant la MGF, montrer que si $X_1 \sim \mathcal{N}(\mu_1, \sigma_1^2)$ et
$X_2 \sim \mathcal{N}(\mu_2, \sigma_2^2)$ sont ind\'ependantes, alors
$X_1 + X_2 \sim \mathcal{N}(\mu_1 + \mu_2, \sigma_1^2 + \sigma_2^2)$.
\end{exercice}
